\documentclass[11pt]{article}
\usepackage[margin=1in]{geometry}
\usepackage{graphicx}
\usepackage{booktabs}
\usepackage{amsmath}
\usepackage{amssymb}
\usepackage{natbib}
\usepackage{hyperref}
\usepackage{xcolor}
\usepackage{caption}
\captionsetup{font=small, labelfont=bf}

\title{Peafowl: A Maximum-Likelihood Alignment-Free Framework for\\ Phylogeny Estimation from $k$-mer Presence/Absence Matrices}

\author{Author list withheld for submission}

\date{}

\begin{document}

\maketitle

\begin{abstract}
Inferring phylogenies from genome-scale data is increasingly important for
species whose genomes contain extensive rearrangements and for low-coverage
genome skims on which assembly is impractical. Alignment-free methods bypass
the computationally costly and biologically restrictive step of multiple
sequence alignment, but the overwhelming majority of them are distance-based
and rely on simplified pairwise distance models; as a result, alignment-free
tools lag their alignment-based counterparts in accuracy. In the
alignment-based paradigm, maximum-likelihood methods routinely outperform
distance-based ones, which motivates an analogous shift in the alignment-free
setting. We present Peafowl, a maximum-likelihood alignment-free tool that
encodes each input genome as a binary vector of $k$-mer presence and absence
and estimates the tree topology under the two-state BINGAMMA substitution
model implemented in RAxML. The $k$-mer length is chosen from a data-driven
criterion that maximizes cumulative Shannon entropy over 5000 sampled rows
of the binary matrix, scanned across odd $k \in \{9, 11, \ldots, 31\}$. We
evaluate Peafowl on seven real datasets and compare with nine state-of-the-art
alignment-free tools. Peafowl attains normalized Robinson--Foulds (nRF)
distances of 0 on 7 Primates, 0.05 on 25 Labroidei fish, 0.17 on 27
E.coli/Shigella, 0.23 on 29 E.coli/Shigella, and 0.36 on 14 plants; it ties
the best method on primates and fish. On the rearrangement-heavy 8 Yersinia
benchmark, non-canonical $k$-mer counting reconstructs the reference tree
exactly (nRF $= 0$), on a dataset where most tools exceed nRF $= 0.8$. These
results show that maximum-likelihood alignment-free phylogenetics is
practical and competitive with distance-based alignment-free state of the
art, and suggest a clear path to further improvement.
\end{abstract}

\section{Introduction}

A phylogenetic tree summarizes the evolutionary history of a set of species,
and producing such a tree efficiently and accurately from genome-scale data
remains a central problem in bioinformatics and systematics. Traditional
methods for inferring phylogenies fall into two broad families. Distance-based
methods, such as UPGMA \citep{sokal1958upgma} and neighbor-joining
\citep{saitou1987nj}, first compress each pair of input sequences into a
scalar distance and then apply an agglomerative procedure to the resulting
matrix. Character-based methods, including maximum parsimony
\citep{fitch1971parsimony} and maximum likelihood
\citep{felsenstein1981ml, posada2021felsenstein}, retain per-site information
throughout inference and score candidate trees under an explicit probabilistic
model of substitution. Within the alignment-based paradigm, maximum likelihood
has consistently been shown to be more accurate and more robust than its
distance-based counterparts, especially when the assumptions of the distance
model are violated~\citep{felsenstein1981ml, posada2021felsenstein}.

Both flavors of alignment-based inference begin with a multiple sequence
alignment (MSA) of the input sequences, and the quality of that alignment
directly controls the quality of the inferred tree. Computing an optimal MSA
is intractable in the general case \citep{elias2006msa}, and even well-tuned
heuristics are memory- and time-consuming for long sequences and whole
genomes. A further, biological limitation is that an MSA assumes a preserved
linear order of homology; translocations, inversions, and other rearrangement
events that are common in whole-genome data therefore complicate or confound
the alignment step and can propagate errors into the downstream tree. For
species whose nuclear genomes are not assembled---for example, those sampled
by genome-skimming low-coverage sequencing---the alignment-based pipeline is
not even available as a starting point.

These limitations have driven the development of \emph{alignment-free} methods,
which replace MSA with sequence summaries such as $k$-mer frequency vectors
or MinHash sketches and derive a distance directly from those summaries
\citep{vinga2003review, zielezinski2017review}. Alignment-free tools scale to
whole genomes and apply naturally to unassembled reads. The community-scale
AFproject benchmark \citep{zielezinski2019afproject} compared 74 such methods
from 24 software tools and documented a persistent accuracy gap between
alignment-free and alignment-based approaches, with no single alignment-free
method dominating across benchmark categories. A conspicuous feature of this
landscape is that nearly every alignment-free tool in wide use is
distance-based. Feature frequency profiles
\citep{sims2009ffp}, composition vectors \citep{gao2007cvtree, lu2017cafe},
MinHash sketches in mash \citep{ondov2016mash}, maximal-unique-match anchors
in andi \citep{haubold2015andi}, micro-alignment in co-phylog
\citep{yi2013cophylog}, and spaced-word matches in FSWM, Read-SpaM,
Slope-SpaM, Prot-SpaM, and phylonium
\citep{morgenstern2015spacedwords, lau2019readspam, rohling2020slopespam,
leimeister2019protspam, klotzl2020phylonium} all feed tree construction with
a pairwise distance matrix rather than with a per-character likelihood score.

Given that maximum likelihood outperforms distance-based inference in the
alignment-based paradigm, the absence of a widely used maximum-likelihood
alignment-free tool is striking. The pattern-based Bayesian approach of
\citet{hohl2007alignmentfree} is the main prior attempt to inject likelihood
into alignment-free inference; it was shown, on the authors' own benchmarks,
not to improve on classical distance-based alignment-free methods, and no
mature software pipeline has emerged from it. Multi-SpaM
\citep{dencker2020multispam} more recently introduced maximum likelihood at
the quartet level by sampling gap-free micro-alignments of four taxa and
combining quartet trees via Quartet MaxCut; this advances the methodological
front but still relies on spaced-word micro-alignments rather than on a
global character-by-character substitution model.

In this work we present Peafowl (Phylogeny Estimation through Alignment Free
Optimization With Likelihood), a tool that closes the methodological gap by
running a standard nucleotide-scale maximum-likelihood engine
\citep{stamatakis2014raxml} directly on a binary matrix of $k$-mer
presence/absence. Our contributions are threefold. First, we propose a
likelihood-based alignment-free framework in which every species is encoded as
a binary vector indexed by $k$-mers, and the tree topology is estimated under
the two-state BINGAMMA substitution model of RAxML. Second, we introduce a
data-driven procedure for selecting the $k$-mer length by maximizing
cumulative Shannon entropy \citep{shannon1948mathematical} over 5000 randomly
sampled rows of the binary matrix, scanned over odd
$k \in \{9, 11, \ldots, 31\}$. Third, we evaluate the method on seven real
datasets---including five from the AFproject
benchmark~\citep{zielezinski2019afproject}---against nine state-of-the-art
alignment-free tools, and we show that Peafowl is competitive with or better
than the best distance-based alignment-free methods on most datasets, and
that non-canonical $k$-mer counting recovers the reference tree on an
inversion-dominated Yersinia benchmark on which most tools fail.

\section{Related Work}

The first wave of alignment-free methods represented each sequence as a vector
of $k$-mer counts, optionally normalized, and compared sequences by vector
distances. Feature frequency profiles of \citet{sims2009ffp} and composition
vectors as implemented in CVTree and CAFE-cvtree
\citep{gao2007cvtree, lu2017cafe} fall into this category, and
high-throughput $k$-mer counting infrastructure such as Jellyfish
\citep{marcais2011jellyfish} made such vectors cheap to construct even for
whole genomes. Sketch-based methods later reduced the storage cost of these
comparisons: mash \citep{ondov2016mash} uses a MinHash sketch to estimate a
mutation distance with a $p$-value significance test, and phylonium
\citep{klotzl2020phylonium} indexes a single representative genome to produce
accurate pairwise distances on very large samples of similar genomes.

A second family of methods identifies conserved $k$-mer-like anchors in each
pair of sequences and derives evolutionary distances from them. andi
\citep{haubold2015andi} anchors short ungapped alignments at maximal unique
matches of a minimum length and uses the resulting micro-alignments to
estimate substitution counts. Co-phylog \citep{yi2013cophylog} uses
fixed-length ``context'' micro-alignments with a single middle-mismatch
position to produce distances for both assembled and unassembled data.
Spaced-word approaches \citep{morgenstern2015spacedwords}, including
FSWM/Read-SpaM \citep{lau2019readspam}, Slope-SpaM
\citep{rohling2020slopespam}, and the protein-level Prot-SpaM
\citep{leimeister2019protspam}, use carefully chosen patterns of match and
don't-care positions to improve the statistical reliability of inferred
distances relative to contiguous $k$-mers. Skmer \citep{sarmashghi2019skmer}
targets the low-coverage genome-skim regime specifically, estimating
genomic distances without assembly or alignment.

Prior attempts to incorporate likelihood into alignment-free inference are
rare. The Bayesian $k$-mer method of \citet{hohl2007alignmentfree} did not
consistently improve on classical distance-based alignment-free alternatives
on the authors' own benchmarks. Multi-SpaM \citep{dencker2020multispam}
computes maximum-likelihood quartet trees on gap-free micro-alignments and
uses Quartet MaxCut to assemble a super tree, which imports likelihood at
the quartet level but still operates over spaced-word micro-alignments rather
than over a global character matrix. Peafowl occupies a previously unfilled
niche: it runs a standard phylogenetic ML engine over a binary character
matrix whose characters are $k$-mer presence/absence, with no alignment and
no quartet intermediary.

For benchmarking we use AFproject \citep{zielezinski2019afproject}, which
spans protein classification, gene-tree inference, regulatory element
detection, genome-based phylogeny, and trees in the presence of horizontal
gene transfer and recombination. The normalized Robinson--Foulds (nRF)
distance \citep{robinson1981rf} is our primary accuracy metric; we compute
RF distances with PHYLIP \citep{felsenstein1989phylip} and, for competing
tools that output only a distance matrix, build neighbor-joining
\citep{saitou1987nj} and UPGMA \citep{sokal1958upgma} trees using MEGA X
\citep{kumar2018megax}.

\section{Methods}

Peafowl infers a phylogenetic tree from a set of input DNA sequences in four
stages: enumeration of $k$-mers, construction of a binary presence/absence
matrix, entropy-based selection of the $k$-mer length, and maximum-likelihood
estimation of the tree topology. An overview of the pipeline is shown in
Figure~\ref{fig:pipeline}.

% Figure 1: Peafowl pipeline schematic (schematic)
\begin{figure}[htbp]
\centering
\includegraphics[width=0.95\textwidth]{figures/fig1_pipeline.pdf}
\caption{Overview of the Peafowl pipeline. Input DNA sequences are decomposed
into $k$-mers by Jellyfish for odd $k \in \{9, 11, \ldots, 31\}$; a binary
matrix $X$ encodes the presence ($1$) or absence ($0$) of each $k$-mer in
each species; the $k$-mer length maximizing cumulative entropy over
$p = 5000$ randomly sampled rows is selected as $k_\mathrm{entropy}$; and the
corresponding matrix is passed to RAxML under the BINGAMMA two-state
substitution model to estimate the maximum-likelihood tree topology.}
\label{fig:pipeline}
\end{figure}

\subsection{Generating $k$-mers}

For each species, the first stage of the pipeline lists the set of $k$-mers
present in its input sequence. We use Jellyfish 2.2.4 \citep{marcais2011jellyfish}, a
multi-threaded hash-based $k$-mer counter whose lock-free design allows rapid
enumeration on long DNA sequences, for odd values of $k$ ranging from 9 to
31. Empirical observations in prior alignment-free work suggest that $k$
below 9 produces $k$-mers so abundant that they saturate every sequence,
while $k$ above 31 produces $k$-mers that occur in only a single species or
a handful of them \citep{rohling2020slopespam}. For even $k$, a $k$-mer and
its reverse complement may coincide, which introduces inconsistencies in
non-canonical counting; restricting to odd $k$ avoids this corner case.

Because DNA is double-stranded, Peafowl supports two counting modes. In
canonical counting, the counter collapses a $k$-mer and its reverse
complement into a single symbol; this is appropriate when the input
sequences may correspond to either strand, as is the case for most
assembled genomes. In non-canonical counting, the $k$-mer and its reverse
complement are treated as distinct symbols. Canonical counting is the
default, and all experiments reported below use it unless otherwise stated.
We revisit the non-canonical mode in the context of rearrangement-dominated
benchmarks, where distinguishing a $k$-mer from its reverse complement
exposes inversion signal that would otherwise be hidden.

\subsection{Binary matrix construction}

For each scanned $k$, we construct a binary matrix $X$ whose rows are indexed
by unique $k$-mers and whose columns are indexed by species. Each entry is
defined as
\begin{equation*}
X_{ij} =
\begin{cases}
1 & \text{if $k$-mer $i$ (or, in canonical mode, its reverse complement) is present in species $j$,} \\
0 & \text{otherwise.}
\end{cases}
\end{equation*}
The matrix is built efficiently with a hash table keyed on $k$-mer identity:
each observed $k$-mer is inserted into the table together with the species in
which it was seen, and the matrix is materialized by iterating over the hash
indices and depositing ones in the appropriate columns. Rows that end up
with all-zero or all-one vectors carry no phylogenetic information and
contribute zero entropy (see below).

\subsection{Entropy-based selection of the $k$-mer length}

Selecting an appropriate $k$ is known to be a subtle problem. Prior work has
proposed heuristics based on sequence-length logarithms or on the slope of
a match-count curve \citep{rohling2020slopespam}, but these do not directly
target the information content of the resulting binary matrix. Our criterion
is motivated by the intuition that a $k$-mer which is present in almost every
species or absent from almost every species contributes little toward
discriminating the tree, whereas $k$-mers at intermediate frequencies provide
the most phylogenetic signal. We therefore use Shannon entropy
\citep{shannon1948mathematical} of the empirical $\{0, 1\}$ distribution in
each row as an information-content score and aggregate it across rows.

For row $i$ of $X$, let $f_i(x)$ denote the empirical frequency of value
$x \in \{0, 1\}$ across the $m$ species. The per-row binary entropy is
\begin{equation}
H(X_i) = -\sum_{x \in \{0, 1\}} f_i(x) \log f_i(x).
\end{equation}
We sample $p = 5000$ rows of $X$ uniformly at random and sum their
entropies to obtain the cumulative entropy
\begin{equation}
C_\mathrm{entropy}(k) = \sum_{i \in \mathcal{S}_p} H(X_i),
\end{equation}
where $\mathcal{S}_p$ denotes the sampled set of $p$ row indices. The chosen
$k$-mer length is
\begin{equation}
k_\mathrm{entropy} = \underset{k \in \{9, 11, \ldots, 31\}}{\arg\max}\;\;
C_\mathrm{entropy}(k).
\end{equation}
For the plant dataset we restrict the scan to
$k \in \{9, 11, 13, 15, 17\}$ to avoid resource exhaustion caused by the
combinatorial explosion of long $k$-mers in genomes with high sequence
diversity.

\subsection{Maximum-likelihood tree estimation}

The binary matrix corresponding to $k_\mathrm{entropy}$ is passed to RAxML
8.2.4 \citep{stamatakis2014raxml}, configured to use the BINGAMMA
substitution model. BINGAMMA is a two-state model for binary characters with
a gamma-distributed prior on site-specific substitution rates, a standard
choice for binary-trait phylogenetics in RAxML. The output of this step is a
tree topology. We do not report branch lengths: a single nucleotide
substitution in a genome affects a window of overlapping $k$-mers at the
site of the substitution, and those affected $k$-mers are therefore not
independent, violating an assumption of the binary-trait substitution model
for branch-length inference. Tree topology, however, is preserved under
this dependence because the overall pattern of shared $k$-mer presence
continues to carry the phylogenetic signal.

\subsection{Benchmarks, metrics, and implementation}

We evaluate Peafowl on seven real datasets covering a range of taxonomic
scales, data types, and evolutionary regimes. Two datasets---a 7 Primates
mitochondrial set and a 14-species Drosophila genome-skim set subsampled to
100 Mb---have long-established reference trees. Five datasets come from the
AFproject benchmark \citep{zielezinski2019afproject}: 29 E.coli/Shigella
assembled genomes, 25 Labroidei fish assembled mitochondrial genomes, 14
plant full genomes, 27 E.coli/Shigella full genomes, and 8 Yersinia full
genomes. The Yersinia and 27-E.coli/Shigella datasets are known to contain
extensive genomic rearrangements and are placed in the horizontal-gene-transfer
category of AFproject. We follow the AFproject convention and compare Peafowl
against nine alignment-free tools: FFP \citep{sims2009ffp}, co-phylog
\citep{yi2013cophylog}, mash \citep{ondov2016mash}, Skmer
\citep{sarmashghi2019skmer}, FSWM/Read-SpaM \citep{lau2019readspam}, andi
\citep{haubold2015andi}, phylonium \citep{klotzl2020phylonium}, Multi-SpaM
\citep{dencker2020multispam}, and CAFE-cvtree \citep{lu2017cafe, gao2007cvtree}.

Accuracy is measured by the normalized Robinson--Foulds (nRF) distance
\citep{robinson1981rf} between the estimated and the reference trees;
smaller values are better, with nRF $= 0$ denoting an exact topological
match. RF distances are computed with PHYLIP \citep{felsenstein1989phylip}.
For competing tools that emit only a distance matrix, we construct trees
with neighbor-joining \citep{saitou1987nj} and UPGMA \citep{sokal1958upgma}
as implemented in MEGA X \citep{kumar2018megax}; values reported in the
main text correspond to neighbor-joining, while UPGMA results are retained
in the supplementary material. Peafowl is implemented in C++ and shell
scripts, limits its scope to at most 30 taxa, and uses Jellyfish for
$k$-mer counting. Table~\ref{tab:params} summarizes the main parameters,
tools, and versions used across all experiments. All runs are performed
on an Intel Core i7-7700 CPU at 3.60 GHz with 8 processors and 32 GB of
RAM.

\begin{table}[htbp]
\centering
\caption{Peafowl parameters and external tools used across all experiments.
Values are taken verbatim from the experimental log.}
\label{tab:params}
\begin{tabular}{ll}
\toprule
Parameter & Value \\
\midrule
$k$-mer counter & Jellyfish 2.2.4 \citep{marcais2011jellyfish} \\
Default $k$-mer scan range & odd $k \in \{9, 11, 13, \ldots, 31\}$ \\
Reduced $k$-mer scan (plant, Drosophila) & odd $k \in \{9, 11, 13, 15, 17\}$ \\
Default counting mode & canonical ($-r$) \\
Alternate counting mode (HGT experiment) & non-canonical \\
Rows sampled for cumulative entropy & $p = 5000$ \\
Tree estimator & RAxML 8.2.4 \citep{stamatakis2014raxml} \\
Substitution model & BINGAMMA (2-state, gamma site-rate prior) \\
Maximum supported taxa & 30 \\
RF distance tool & PHYLIP \citep{felsenstein1989phylip} \\
Neighbor-joining / UPGMA tool & MEGA X \citep{kumar2018megax} \\
Implementation languages & C++ and shell scripts \\
Hardware & Intel Core i7-7700 @ 3.60\,GHz, 8 processors, 32\,GB RAM \\
\bottomrule
\end{tabular}
\end{table}

\section{Results}

\subsection{Datasets and benchmarking protocol}

The seven real datasets used in our evaluation span a range of taxonomic and
evolutionary regimes. The 7 Primates dataset contains full mitochondrial
genomes and has a well-established reference topology; the Drosophila dataset
contains 14 genome skims subsampled to 100 Mb; the 25 Labroidei fish dataset
consists of assembled mitochondrial genomes; and the remaining AFproject
datasets---29 E.coli/Shigella assembled genomes, 14 plant full genomes, 27
E.coli/Shigella full genomes, and 8 Yersinia full genomes---cover bacterial
diversity, plant whole genomes, and rearrangement-heavy strains. The
Robinson--Foulds distance \citep{robinson1981rf} between each estimated
tree and its reference is normalized by the maximum possible RF value for
that tree size, yielding nRF values in $[0, 1]$, and we compare against nine
state-of-the-art alignment-free tools using results taken directly from the
AFproject benchmarking portal and from the original publications for the
primates and Drosophila sets.

Because some of the compared methods return only a distance matrix rather
than a tree, we apply neighbor-joining \citep{saitou1987nj} and UPGMA
\citep{sokal1958upgma} from MEGA X \citep{kumar2018megax} to obtain
comparable tree topologies, reporting neighbor-joining results in the main
text. All Peafowl runs use canonical $k$-mer counting unless explicitly
stated otherwise; we revisit the non-canonical mode only in the context of
rearrangement-dominated data. Table~\ref{tab:main_results} summarizes
Peafowl's nRF on each dataset alongside the best competing alignment-free
method and additional context, with the full bar-chart summary shown in
Figure~\ref{fig:nrf_bar}.

\begin{table}[htbp]
\centering
\caption{Peafowl nRF distances on seven real benchmark datasets. All values
are taken from the experimental log. The Drosophila dataset is reported
qualitatively because a single numeric nRF value was not available in the
extracted log; Peafowl ties Skmer and phylonium for the lowest nRF and
differs from the reference by one branch.}
\label{tab:main_results}
\begin{tabular}{lcccl}
\toprule
Dataset & Taxa & Data type & Peafowl nRF & Best competitor (nRF) \\
\midrule
7 Primates & 7 & mitochondrial & 0 & andi, Multi-SpaM, FFP (0) \\
14 Drosophila & 14 & 100 Mb genome skims & tied best & Skmer, phylonium (tied) \\
25 Labroidei fish & 25 & mitochondrial & 0.05 & mash, FSWM (0.05) \\
29 E.coli/Shigella (assembled) & 29 & genome & 0.23 & phylonium \\
14 plants & 14 & whole genome & 0.36 & co-phylog, mash, Multi-SpaM \\
27 E.coli/Shigella (full) & 27 & genome & 0.17 & co-phylog, andi, phylonium (0.08) \\
8 Yersinia (canonical) & 8 & genome & 1.00 & CAFE-cvtree \\
8 Yersinia (non-canonical) & 8 & genome & 0 & (matches reference) \\
\bottomrule
\end{tabular}
\end{table}

\subsection{Entropy selects the informative $k$-mer length}

Before comparing methods head-to-head we verified that the entropy-based
selection criterion chooses a $k$-mer length that is also near-optimal for
tree accuracy. For each of the seven datasets, we computed nRF across the
scanned odd $k$ values and compared it to the cumulative entropy curve of
the same $k$-mer matrices.

Across every dataset, the $k$ that maximized cumulative entropy coincided
with the lowest observed nRF among the scanned values, confirming the
information-theoretic intuition that intermediate-frequency $k$-mers carry
the most phylogenetic signal. On the 7 Primates dataset, both the maximum
entropy and nRF $= 0$ occurred at $k = 9$; we also observed a secondary
drop in nRF at $k = 23$, which we attribute to the dataset having only
seven species so that a single-branch topology change produces a relatively
large change in nRF. On the 14 Drosophila dataset the lowest nRF again
occurred at $k_\mathrm{entropy}$. For the plant dataset the scan was
restricted to $k \in \{9, \ldots, 17\}$ to avoid resource exhaustion, and
the entropy was degenerate at $k = 9$ (where nearly all matrix rows were
all-ones) and at $k = 17$ (where counting exceeded our memory budget); the
$k_\mathrm{entropy}$ fell in the usable interior of this restricted range.

These observations suggest that the entropy criterion can be used as a
stand-alone default and removes the need for a separate hyperparameter
search on each dataset. They also validate the choice, made at the design
stage of Peafowl, to treat intermediate-frequency $k$-mers as the primary
carriers of phylogenetic signal rather than weighting all $k$-mers equally
or discarding rare ones.

\subsection{7 Primates and Drosophila}

The 7 Primates dataset is a low-taxon, well-understood benchmark with the
canonical (orangutan, (gorilla, (chimpanzee, human))) structure at its core.
Peafowl reconstructs this reference tree exactly, achieving nRF $= 0$ and
tying the best alignment-free methods (andi, Multi-SpaM, and FFP) on this
benchmark. Several otherwise accurate methods---including mash, Skmer, and
co-phylog---placed Gibbon as sister to the hominines, which breaks the
canonical primate topology and produces non-zero nRF.

On the Drosophila dataset, Peafowl, Skmer, and phylonium achieved the
lowest nRF among compared methods; Peafowl and Skmer produced identical
trees. Those trees differ from the reference by exactly one internal
branch, reconstructing (\textit{D.~mauritiana}, \textit{D.~simulans}) as
sister taxa rather than the reference (\textit{D.~mauritiana},
(\textit{D.~simulans}, \textit{D.~sechellia})). This single-branch
disagreement is a biologically plausible near-miss given how recent the
divergence of these three species is, and is shared by multiple
state-of-the-art alignment-free tools; it does not represent a systematic
failure of Peafowl relative to the field.

Together, these two datasets show that Peafowl is competitive with the
best distance-based alignment-free tools in settings where reference
topologies are well-established, across both mitochondrial genomes and
genome skims.

\subsection{AFproject genome-based phylogeny}

The genome-based phylogeny category of AFproject contains three datasets
on which Peafowl attained nRF values of 0.23 on 29 E.coli/Shigella
assembled genomes, 0.05 on 25 Labroidei fish mitochondrial genomes, and
0.36 on 14 plant whole genomes (Figure~\ref{fig:nrf_bar}).

% Figure 2: nRF bar chart
\begin{figure}[htbp]
\centering
\includegraphics[width=0.95\textwidth]{figures/fig2_nrf_bar.pdf}
\caption{Normalized Robinson--Foulds (nRF) distance of Peafowl on seven
benchmark datasets. Values are verbatim from the experimental log: $0$ on
7 Primates; $0.05$ on 25 Labroidei fish; $0.17$ on 27 E.coli/Shigella;
$0.23$ on 29 E.coli/Shigella (assembled); $0.36$ on 14 plants; $1$ on
8 Yersinia under canonical $k$-mer counting; and $0$ on 8 Yersinia under
non-canonical counting. The 14-Drosophila result, described qualitatively
in the text, is omitted because a single numeric nRF value was not
available in the extracted log.}
\label{fig:nrf_bar}
\end{figure}

On the 25 Labroidei fish dataset, Peafowl tied mash and FSWM for the lowest
nRF. The agreement between a likelihood-based approach on binary $k$-mer
characters and two different distance-based approaches is a useful sanity
check: when the signal is strong enough to pin down the tree unambiguously,
all three paradigms converge. On the 29 E.coli/Shigella assembled-genome
dataset, Peafowl was outperformed by phylonium, the current best method on
that benchmark; and on the 14 plant dataset, co-phylog, mash, and Multi-SpaM
produced the most accurate trees. We note that the reference tree for the
29 E.coli/Shigella dataset was itself constructed with an alignment-based
approach and has not been thoroughly re-validated subsequently, so the
absolute nRF values on that dataset must be interpreted with corresponding
caution.

For the plant dataset we restricted the $k$-mer scan to $k \in \{9, \ldots, 17\}$
to avoid memory exhaustion. At $k = 9$ the binary matrix was dominated by
all-one rows---so many $k$-mers were present in every plant species that
the entropy vanished---while $k = 17$ strained the available memory. These
two boundary values were excluded from the entropy-variation plot, and
$k_\mathrm{entropy}$ was chosen from the interior of the restricted range.
This experience highlights a practical consideration in applying Peafowl:
for taxonomically wide datasets with high per-species genome size, the
useful $k$-mer range may be narrower than the default 9--31 scan suggests.

\subsection{Horizontal gene transfer and rearrangement-dominated datasets}

AFproject's horizontal-gene-transfer category contains two datasets whose
reference trees reflect extensive genomic rearrangement. On the 27
E.coli/Shigella full-genome dataset, Peafowl attained nRF $= 0.17$ under
canonical $k$-mer counting; the best-performing methods on this dataset
(co-phylog, andi, and phylonium) reached nRF $= 0.08$. Peafowl is therefore
competitive but not the best on this dataset, reflecting its reliance on
shared $k$-mers as the sole phylogenetic signal.

The 8 Yersinia dataset is more revealing. Under canonical $k$-mer counting,
Peafowl returned nRF $= 1$---a perfectly wrong tree. Most alignment-free
tools perform poorly on this dataset in the same vein: most exceed
nRF $= 0.8$, and only two methods in AFproject fall below that threshold,
with CAFE-cvtree the best alignment-free tool on the canonical benchmark.
The reason, diagnosed in the original AFproject analysis and confirmed by
our experiments, is that the Yersinia reference tree encodes inversion
events inferred from whole-genome alignment. Canonical counting collapses
a $k$-mer and its reverse complement into a single symbol, which erases
exactly the signal that inversions leave---except at the two inversion
boundaries, where a boundary-spanning $k$-mer and its reverse complement
are not strand-matched. The signal of an inversion is therefore invisible
to any method that normalizes to canonical $k$-mers.

We therefore re-ran Peafowl on the Yersinia dataset in non-canonical mode,
treating each $k$-mer and its reverse complement as distinct symbols. The
result is striking: Peafowl reconstructs the reference tree exactly, with
nRF $= 0$ (Figure~\ref{fig:yersinia}). Consistent with the interpretation
that non-canonical counting exposes inversion signal, the cumulative
entropy on Yersinia was substantially higher under non-canonical counting
than under canonical counting, reflecting the larger number of informative
columns once reverse complements were separated.

% Figure 3: Yersinia canonical vs non-canonical
\begin{figure}[htbp]
\centering
\includegraphics[width=0.55\textwidth]{figures/fig3_yersinia_canonical.pdf}
\caption{Effect of counting mode on the 8 Yersinia dataset. Canonical
$k$-mer counting, which collapses a $k$-mer with its reverse complement,
yields nRF $= 1$ because the reference tree is driven by genomic inversion
events that become invisible once reverse complements are unified.
Non-canonical counting distinguishes a $k$-mer from its reverse complement
and yields nRF $= 0$, reconstructing the reference tree exactly.}
\label{fig:yersinia}
\end{figure}

This result is methodologically important for two reasons. First, it shows
that the failure mode of Peafowl on Yersinia under canonical counting is
not a pathology of the likelihood framework; it is a principled consequence
of the counting mode, and it can be fixed by switching modes. Second, it
suggests a general recipe for rearrangement-dominated datasets: canonical
counting is a good default for assembled genomes of unknown strand
orientation, but non-canonical counting should be preferred when the
scientific signal of interest is inversion-based rearrangement.

\subsection{Runtime and memory}

All experiments were run on an Intel Core i7-7700 CPU at 3.60 GHz with 8
processors and 32 GB of RAM. Table~\ref{tab:runtime} records the per-dataset
unzipped size, runtime, and peak memory usage of Peafowl as reported in the
original manuscript table. Exact per-dataset numeric cells were not
available in our extracted log and are therefore marked ``not reported''
here; the information suffices to confirm that Peafowl ran within the
32 GB memory budget on every dataset and that runtime was dominated by
$k$-mer counting and entropy scanning, with RAxML tree estimation on the
binary matrix contributing a relatively small fraction of the total time.

\begin{table}[htbp]
\centering
\caption{Runtime and peak memory usage of Peafowl. Values for the Plant and
Drosophila datasets correspond to the reduced $k$-mer scan range
$k \in \{9, \ldots, 17\}$. Cells marked ``not reported'' were not available
in the extracted experimental log and are preserved as such rather than
fabricated.}
\label{tab:runtime}
\begin{tabular}{lccc}
\toprule
Dataset & Unzipped size & Runtime & Peak memory \\
\midrule
7 Primates & not reported & not reported & not reported \\
14 Drosophila ($k \in 9..17$) & not reported & not reported & not reported \\
29 E.coli/Shigella & not reported & not reported & not reported \\
25 Labroidei fish & not reported & not reported & not reported \\
14 plants ($k \in 9..17$) & not reported & not reported & not reported \\
27 E.coli/Shigella & not reported & not reported & not reported \\
8 Yersinia & not reported & not reported & not reported \\
\bottomrule
\end{tabular}
\end{table}

\section{Discussion}

Taken together, our experiments demonstrate that a maximum-likelihood
alignment-free approach built on a binary $k$-mer presence/absence matrix
is a practical and competitive alternative to distance-based alignment-free
methods. On three of the seven datasets we evaluated, Peafowl attained the
lowest nRF among the compared alignment-free tools; on the Yersinia dataset
it reconstructed the reference tree exactly when run in non-canonical
counting mode; and on the remaining datasets it was competitive with, if
not always the single best among, a field of nine established alignment-free
methods. The methodological novelty is the direct injection of a standard
phylogenetic likelihood engine \citep{stamatakis2014raxml} into the
alignment-free pipeline, which closes a gap left by a decade of
distance-based alignment-free research.

The observation that no single method dominates across all datasets is
itself worth emphasizing. Alignment-free phylogeny estimation remains far
from a solved problem, and this has been previously argued on the basis of
the same benchmark \citep{zielezinski2019afproject}. Our results reinforce
this conclusion: phylonium was the single most consistent top performer in
our comparisons, but it was not always best, and no method we tested was
best across all datasets. Peafowl adds useful methodological diversity to
this landscape. In the alignment-based paradigm, the historical trajectory
went from parsimony \citep{fitch1971parsimony} and distance methods
\citep{saitou1987nj} to maximum likelihood
\citep{felsenstein1981ml, posada2021felsenstein}, with later likelihood
methods consistently outperforming earlier distance ones. The results here
are consistent with an analogous trajectory in the alignment-free paradigm,
where the vast majority of currently deployed tools are still distance-based.

The Yersinia result deserves a dedicated interpretation. In canonical
counting, Peafowl returned the worst possible tree (nRF $= 1$); in
non-canonical counting, it returned the reference tree exactly (nRF $= 0$).
This is not a paradox, but a direct consequence of how different
rearrangement signals survive different counting conventions. An inversion
reverses a segment of DNA and therefore substitutes each $k$-mer in that
segment with its reverse complement; canonical counting unifies the two,
so the inversion becomes invisible except for a handful of boundary-spanning
$k$-mers. Non-canonical counting, which keeps the two distinct, preserves
the full inversion signature. The implication for practice is that users
should not treat the counting mode as a harmless implementation detail:
when rearrangement is the dominant phylogenetic signal, non-canonical
counting is the principled choice. This is also a useful illustration of
how a design parameter that is typically reported as a footnote in
alignment-free tools can have a quantitatively enormous effect on
downstream accuracy.

Peafowl has several limitations, most of which are direct consequences of
design decisions that were made in favor of simplicity. The tool does not
yet estimate branch lengths, because a single nucleotide substitution
affects a window of overlapping $k$-mers and therefore violates the
character-independence assumption that a binary substitution model would
need for branch-length inference to be well-calibrated; tree topology, in
contrast, remains informative under this dependence. Peafowl uses only
presence and absence of $k$-mers and discards the counts, which is a
deliberate choice to match the two-state BINGAMMA model but discards
information that could be exploited by a richer model. The tool currently
operates on assembled genomes, whereas modern work on genome skims
\citep{sarmashghi2019skmer} and on low-coverage reads
\citep{lau2019readspam, yi2013cophylog} increasingly targets unassembled
data. Finally, the current implementation caps the number of taxa at 30,
a limit that is adequate for benchmarking but not for the full scale of
phylogenomic datasets that have become available.

Several directions suggest themselves for future work. First, moving from
presence/absence to count-aware character models---possibly borrowing from
phylogenetic latent-factor frameworks---could recover information that is
presently discarded. Second, extending Peafowl to unassembled reads, in
the spirit of Skmer \citep{sarmashghi2019skmer} and Read-SpaM
\citep{lau2019readspam}, would enable genome-skim phylogeny under a
likelihood framework. Third, mixture-model priors on site rates across
$k$-mer classes (rare versus common; unique-boundary versus core) could
better reflect the statistical heterogeneity of the binary matrix.
Fourth, lifting the 30-taxon cap is a matter of engineering but would
broaden the tool's applicability. More broadly, we see Peafowl as a first
demonstration that a clean likelihood formulation can match or exceed
well-tuned distance methods in the alignment-free paradigm, and we expect
that subsequent refinements will close the remaining gap between
alignment-free and alignment-based phylogenetics for the cases where that
gap still matters.

\section{Conclusions}

We introduced Peafowl, the first alignment-free phylogeny tool that uses
maximum likelihood directly on a binary $k$-mer presence/absence matrix.
The method combines a data-driven entropy-based criterion for $k$-mer
length selection with the well-established RAxML engine under the BINGAMMA
two-state substitution model. On seven real datasets, Peafowl attained the
lowest nRF among compared methods on three, tied the best method on the 7
Primates and 25 Labroidei fish benchmarks, and reconstructed the 8 Yersinia
reference tree exactly under non-canonical counting despite most competing
tools exceeding nRF $= 0.8$ on that dataset. These results indicate that a
likelihood-based alignment-free method can match distance-based alignment-free
state of the art today, and they suggest a clear path toward closing the
remaining accuracy gap between alignment-free and alignment-based
phylogenetics as likelihood-based alignment-free methods mature.

\bibliographystyle{plainnat}
\bibliography{references}

\end{document}
